\chapter{Code(s) Available and Used}\label{ch:codes}

\section{Closest existing codes on the World Wide Web}

The guideline requires an explicit answer to the question of which existing
code comes closest to this project. A search over Python, Julia, MATLAB,
Fortran/Abaqus and C++ resources found \emph{no} public code that combines all
four ingredients of this project: (i) a periodic finite-strain RVE with an
exactly rigid, freely rotating inclusion; (ii) a condensed consistent tangent;
(iii) a polyconvex neural energy trained on that data; and (iv) deployment of the
network as a material law inside NGSolve. Each ingredient on its own, however,
has close precedents, summarized in \cref{tab:available}. The closest code to the
neural material law in NGSolve is the forum example~\cite{NGSolveForumTF}; the
closest RVE formulations are the COMET and Ferrite.jl homogenization examples.

{\small
\begin{xltabular}{\textwidth}{@{}p{3.1cm}p{1.9cm}X@{}}
\caption{Closest publicly available codes and their relation to this project.}\label{tab:available}\\
\toprule
Code & Language & Functionality and relation to this project\\
\midrule
\endfirsthead
\caption[]{\emph{(continued)}}\\
\toprule
Code & Language & Functionality and relation to this project\\
\midrule
\endhead
\bottomrule
\endfoot
COMET / comet-fenicsx tours~\cite{CometFenics} & Python (FEniCSx, \code{dolfinx\_mpc}) &
  Periodic homogenization of a 2D plane-strain \emph{linear} elastic composite with
  deformable circular inclusions; the periodic fluctuation is the unknown, as in
  \cref{eq:kin}. Closest RVE formulation, but small strain, two-dimensional, no
  rigid inclusion and no neural surrogate.\\
Ferrite.jl tutorial ``Computational homogenization''~\cite{Ferrite} & Julia &
  2D small-strain RVE with inclusions ten times stiffer than the matrix, Dirichlet and
  periodic boundary conditions. The homogenized stiffness follows from sensitivity
  fields for unit macroscopic strains with one reused factorization, the linear
  counterpart of the condensed tangent \cref{eq:schur}.\\
Andreassen \& Andreasen~\cite{Andreassen2014} & MATLAB &
  Short educational code for the homogenized elasticity tensor of periodic
  two-phase composites (small strain, voxel mesh).\\
EasyPBC~\cite{Omairey2019} & Python plug-in for Abaqus/CAE &
  Automatic periodic boundary conditions and effective elastic constants of
  user RVEs in Abaqus (linear elasticity).\\
MonolithFE$^2$~\cite{MonolithFE2} & Fortran (Abaqus UMAT/UEL) &
  Open-source FE$^2$ toolbox from TU Bergakademie Freiberg with monolithic and
  staggered algorithms and periodic micro boundary conditions. It solves the micro
  problem concurrently instead of using an offline surrogate, and is the reference
  point for the concurrent alternative discussed in \cref{ch:intro}.\\
CANN~\cite{Linka2021} & Python (TensorFlow) &
  Constitutive artificial neural networks for hyperelasticity trained on stress
  data; invariant-based, physics-structured, but not coupled to an FE solver in
  the public code and not trained on RVE data.\\
Physics-augmented NNs~\cite{Klein2022,Linden2023,Kalina2022} & method papers &
  Polyconvex invariant-based ICNN energies with stress normalization; the
  methodological basis of the pICNN (features, non-negative weights,
  normalization). Not a reusable package for NGSolve.\\
NGSolve i-tutorial 10.2 ``NGSolve with tensor-flow''~\cite{NGSolveTF} & Python (NGSolve + Keras/TensorFlow) &
  The precedent named by the committee. NGSolve solves a Poisson problem on a
  domain deformed by two parameters; a Keras network is trained to map the
  parameters to the nodal solution and is compared with a new NGSolve solve. Same
  offline-surrogate idea (FE data generation, then training), but the network
  replaces the whole solution and is not evaluated inside the FE solver.\\
NGSolve forum example~\cite{NGSolveForumTF} & Python (NGSolve + TensorFlow) &
  \emph{Closest code to the material-law coupling}: a trained network for the
  strain energy as a function of the Green--Lagrange strain is rebuilt from its
  weights as NGSolve coefficient functions and used in a nonlinear elasticity
  solve, with stresses from symbolic differentiation. The reported nonzero energy
  in the undeformed state is the defect removed by the normalization
  \cref{eq:norm}. No PyTorch, no quadrature-point transfer, no RVE data.\\
DNN in sequential multiscale FE~\cite{Troeger2023} & Python surrogate in Fortran FE &
  A purely data-based DNN surrogate of the RVE response (tangent by reverse-mode
  AD), integrated into a Fortran FE code and used in load-step-controlled simulations
  (2D plane strain, small strain). Reported speed-ups over FE$^2$ are 524$\times$
  (TensorFlow) and 6048$\times$ (JAX with just-in-time compilation) for a plate with a hole;
  these are not comparable with the present figures (other problem, hardware and
  reference). Same offline-surrogate idea; different strain regime, no physics
  constraints, other software.\\
\end{xltabular}
}

\section{Code(s) used in original or adapted form}

No source file of the codes in \cref{tab:available} was copied. The following
parts were used or adapted, as declared in the source documentation.

{\small
\begin{xltabular}{\textwidth}{@{}p{3.4cm}p{2.2cm}X@{}}
\caption{Software used in original or adapted form.}\label{tab:adapted}\\
\toprule
Source & Form & Use in this project\\
\midrule
\endfirsthead
\caption[]{\emph{(continued)}}\\
\toprule
Source & Form & Use in this project\\
\midrule
\endhead
\bottomrule
\endfoot
NGSolve periodicity tutorial~\cite{NGSolvePeriodic} & adapted &
  Face identification before meshing and the \code{Periodic} space wrapper
  (\code{src\slash{}rve.py}, \code{src\slash{}mixed\_matrix.py}).\\
NGSolve API (\code{IntegrationRuleSpace}, \code{NumberSpace}, symbolic forms) & original library &
  Quadrature-point data transfer, global macroscopic coordinates, exact
  linearization of the RVE residual.\\
ICNN architecture~\cite{Amos2017} & re-implemented &
  Non-negative hidden weights through softplus and input passthrough at every
  layer (\code{src\slash{}picnn.py}).\\
Invariant features and normalization~\cite{Linden2023,Klein2022} & adapted &
  Feature set $[I_1,I_2,J,-\ln J]$ and stress-free normalization, here in the
  objective pressure form \cref{eq:norm}.\\
PyTorch (\code{nn.Module}, \code{autograd}, \code{torch.func}, Adam) & original library &
  Network definition, AD stress and tangent, training loops.\\
Netgen OCC geometry kernel & original library & Cube with spherical cavity.\\
\end{xltabular}
}

\section{Original contributions}\label{sec:contrib}

The codes in \cref{tab:available} served as references for individual ideas; none
of their source was copied, and none provides any of the components below. The
project code (\RepositoryURL) comprises about 17\,000 non-blank, non-comment lines of
Python in \code{src/}, \code{tests/} (18 verification scripts) and \code{benchmarks/}
(20 study scripts). \Cref{tab:contrib} lists the components written for this project,
the closest precedent of each and what is new.

{\small
\begin{xltabular}{\textwidth}{@{}>{\raggedright\arraybackslash}p{3.0cm}>{\raggedright\arraybackslash}p{3.1cm}>{\raggedright\arraybackslash}X>{\raggedright\arraybackslash}p{1.5cm}@{}}
\caption[Components developed in this project]{Components developed in this project. Lines: non-blank, non-comment lines of the main module(s).}\label{tab:contrib}\\
\toprule
Component & Module(s), lines & New compared with the closest precedent & Verified by\\
\midrule
\endfirsthead
\caption[]{\emph{(continued)}}\\
\toprule
Component & Module(s), lines & New compared with the closest precedent & Verified by\\
\midrule
\endhead
\bottomrule
\endfoot
PyTorch--NGSolve material bridge &
  \code{ngsolve\_umat.py} (939) &
  The forum example~\cite{NGSolveForumTF} rebuilds a fixed network as symbolic
  coefficient functions. Here the PyTorch model is evaluated batch-wise at all
  quadrature points through \code{IntegrationRuleSpace}; stress and the full
  $9\times9$ AD tangent are returned to NGSolve, with energy line search, secant
  predictor and spectral shift. & V0a--e, V1, V4--V8, M-V0\\
Constrained energy network &
  \code{picnn.py} (171) &
  ICNN~\cite{Amos2017} with invariant features~\cite{Linden2023,Klein2022}; added
  volume-fraction context and the objective pressure normalization \cref{eq:norm}. & V2, V2b, V3, V0d, S-V1\\
Periodic mixed RVE with exactly rigid particle &
  \code{rve.py}, \code{mixed\_matrix.py}, \code{rve\_solver.py} (594) &
  COMET~\cite{CometFenics} and Ferrite.jl~\cite{Ferrite} treat small-strain elastic
  inclusions. Here: finite strain, Taylor--Hood $u$--$p$ matrix, rigid trace
  \cref{eq:rigid} imposed exactly, load continuation. & P2-V1--V4, P2-V8\\
Free particle rotation &
  \code{free\_particle.py} (569) &
  No precedent found: rotation as global unknown on $SO(3)$, condensed rotational
  Newton system from moment equilibrium. & P2-V2, P2-V7\\
Work-conjugate homogenization &
  \code{homogenization.py} (153) &
  Particle-reaction term of \cref{eq:hm}, required because the particle is a
  constraint and not a material phase. & P2-V5\\
Consistent homogenized tangent &
  \code{rve\_tangent.py} (315) &
  Ferrite.jl uses linear sensitivity fields; here the finite-strain Schur complement
  \cref{eq:schur} including the rotational unknowns. & P2-V6, P2-V9\\
Dataset generation and audit &
  \code{rve\_dataset.py}, \code{generate\_rve\_dataset\_parallel.py},
  \code{audit\_rve\_dataset.py}, \code{merge\_rve\_training\_data.py} (1459) &
  Deterministic, checkpointed, parallel generation; splits grouped by deformation
  state; hash-based provenance and audit. & P2-V10\\
Training and fine-tuning &
  \code{train.py}, \code{train\_rve.py}, \code{train\_rve\_shear\_adaptation.py} (1630) &
  Energy--stress--tangent loss with training-split scaling, early stopping, frozen
  checkpoints, shear fine-tuning with data replay. & Sec.~\ref{sec:trainperf}\\
Macroscopic application &
  \code{phase6\_macro\_fe\_demo.py} (347) &
  Surrogate as material law in a 3D finite-strain macro problem, $P_1$ and $P_2$. & M-V0, Sec.~\ref{sec:macro}\\
Verification and benchmarks &
  \code{tests/}, \code{benchmarks/} &
  26 tests (\cref{tab:vmatrix}), micromechanics, timing and holdout studies. & --\\
\end{xltabular}
}

\section{External libraries and versions}

\begin{table}[htbp]
\centering\small
\caption[External software]{External software. ``Cluster'' denotes the environment of all production runs; ``Local'' the clean environment of the final test runs and timing study.}
\label{tab:versions}
\begin{tabular}{@{}llll@{}}
\toprule
Software & Role & Cluster (production) & Local (final tests)\\
\midrule
Python & language & 3.11.7 & 3.13.18\\
NGSolve / Netgen & FEM, geometry, symbolic forms & 6.2.2608 & 6.2.2608\\
PyTorch & network, AD, training & 2.15.0+cpu & 2.15.0 (CPU use)\\
NumPy & arrays & 2.4.6 & 2.5.3\\
Matplotlib & plots & 3.11.0 (Phase 1) & 3.11.2\\
\bottomrule
\end{tabular}
\end{table}

All packages are freely available from PyPI (\code{pip install ngsolve torch
numpy matplotlib}); \code{requirements.txt} in the repository lists minimum
versions (\cref{ch:manual}).
